\section{Distributed Systems Ideas}

This chapter covers the ideas that make the platform correct when machines and
networks fail. They are worth learning well: they appear in databases, job
queues, payment systems and cloud schedulers.

\subsection{Pull, not push}

There are two ways to hand out work. In a \emph{push} design the server picks
a worker and sends it the job. In a \emph{pull} design each worker asks for
work when it is ready. distributetr9 pulls:

\begin{lstlisting}[language=Python]
# agent/agent.py (simplified)
while True:
    status, hb = self.heartbeat()          # "I am alive, here is my state"
    status, payload = self.claim()         # "Any work for me?" (long poll)
    attempt = payload.get("attempt")
    if attempt:
        self.execute(attempt)
\end{lstlisting}

Pulling suits lab PCs well:
\begin{itemize}
\item The server never needs to open a connection \emph{to} a PC. PCs sit
  behind Windows firewalls and NAT; outgoing HTTP from the PC almost always
  works.
\item A busy, broken or rebooting PC simply stops asking. The server does not
  need to track whether a push succeeded.
\item The PC decides when it is ready, for example not while a student is
  gaming on it (the \texttt{gpu\_desktop\_busy} check).
\end{itemize}

\paragraph{Long polling.} Asking every second wastes requests; asking every
minute adds a minute of delay. A \emph{long poll} gets both: the agent asks,
and the server holds the request open, checking once a second, until work
appears or 25 seconds pass. The response comes back the moment a job fits.

\paragraph{Other use cases.} GitHub Actions self-hosted runners, GitLab
runners, Jenkins agents in ``inbound'' mode, Celery and Sidekiq workers pulling
from a queue, and Amazon SQS long polling (\texttt{WaitTimeSeconds}) all pull.

\subsection{Heartbeats}

Every few seconds the agent sends \texttt{POST /agent/heartbeat} with its free
disk space, whether the desktop GPU is busy, and whether it is healthy. The
server records \texttt{last\_heartbeat\_at}. A background worker marks nodes
whose heartbeat is older than three intervals as \texttt{offline}, and the
scheduler skips them. The heartbeat \emph{response} carries commands back:
\texttt{drain}, \texttt{quarantine}, or an available agent update. That gives
the server a way to talk to PCs without ever connecting to them.

\subsection{Leases}

A lease is ownership with an expiry date. When a node claims a job it gets the
attempt for 30 seconds. To keep it, the node must renew before time runs out.

\begin{verbatim}
 t=0s    claim       lease_expires_at = 30s
 t=10s   renew       lease_expires_at = 40s
 t=20s   renew       lease_expires_at = 50s
 t=25s   PC loses power ...
 t=50s   lease expires         -> attempt state: unknown
 t=110s  + disconnect grace    -> failed (lease_lost), job re-queued
\end{verbatim}

Why not just wait for the PC to say it failed? Because a crashed PC cannot say
anything. A lease turns ``silence'' into a decision the server can make on its
own clock. The grace period avoids giving up on a PC that only missed a few
renewals during a Wi-Fi hiccup.

\paragraph{Other use cases.} DHCP leases your laptop an IP address; Kubernetes
leader election uses Lease objects; distributed locks in etcd, ZooKeeper and
Redis (Redlock) are leases; Amazon SQS's ``visibility timeout'' is a lease on a
message.

\subsection{Fencing tokens: rejecting zombies}

Leases create a new problem. Picture this:

\begin{enumerate}
\item PC A claims job 7 and starts training.
\item PC A's network drops for two minutes. Its lease expires.
\item The server gives job 7 to PC B, which finishes and completes it.
\item PC A's network returns. It finishes too and calls \texttt{/complete}.
\end{enumerate}

PC A is a \emph{zombie}: it believes it owns a job it lost long ago. If the
server accepted its result, job 7 would have two results, or the wrong one.

The fix is a \textbf{fencing token}. Each attempt gets a fresh random UUID when
it is created. Every call that changes state must carry it, and the server
compares it with the token stored on the attempt:

\begin{lstlisting}[language=Python]
# control/app/scheduler.py
def _require_token(attempt: JobAttempt, token: uuid.UUID) -> None:
    if attempt.fencing_token != token:
        raise AppError(409, "fenced", "stale fencing token", ...)
\end{lstlisting}

PC A's attempt was failed when its lease expired, and job 7's accepted attempt
is now PC B's. PC A's \texttt{/complete} hits the ``attempt superseded'' check
and gets \texttt{409}. The zombie's work is discarded, never promoted. The
simulated failure exam checks exactly this in its \texttt{network\_drop}
scenario.

\paragraph{Other use cases.} Martin Kleppmann's essay ``How to do distributed
locking'' popularised the term. Apache Kafka uses producer epochs, HDFS uses
generation stamps, and Raft and Paxos use terms and ballot numbers: all are
fencing tokens.

\subsection{Boot sessions}

A PC that reboots gets a brand new agent process, but the same node identity.
On registration the server issues a new \texttt{boot\_session\_id} and fails
every attempt the node still owned (\texttt{on\_node\_boot}). Requests carrying
an old boot session get \texttt{409 stale\_boot\_session}. This catches a case
leases alone would only catch slowly: a PC that rebooted in ten seconds has
certainly lost its running container, so there is no point waiting a full lease.

\subsection{Idempotency: safe to repeat}

Networks lose responses. If an agent sends \texttt{/complete} and the reply
never arrives, it cannot know whether the server processed it. The only safe
move is to send it again, so every state-changing call must be
\textbf{idempotent}: doing it twice has the same effect as doing it once.

\begin{itemize}
\item \texttt{/complete} twice with the same token and exit code returns the
  stored result. With a \emph{different} exit code it returns \texttt{409},
  because that signals a real bug.
\item \texttt{/claim} twice from the same node returns the node's existing
  attempt instead of a second one.
\item \texttt{POST /jobs} accepts an \texttt{Idempotency-Key} header. A browser
  that retries a submit after a timeout gets the same job back instead of a
  duplicate:
\end{itemize}

\begin{lstlisting}[language=Python]
# control/app/routers/portal.py (simplified)
if idem_key:
    existing = await session.scalar(
        select(Job.id).where(Job.user_id == user.id,
                             Job.idempotency_key == idem_key))
    if existing is not None:
        return serialize_job(await load_job(session, existing))
\end{lstlisting}

A unique index on \texttt{(user\_id, idempotency\_key)} is the backstop: if two
identical requests race past the check at the same moment, the second insert
fails and the handler returns the job the first one created.

\paragraph{Other use cases.} Stripe's API is the textbook example: every
payment request takes an \texttt{Idempotency-Key} so a retried charge never
bills twice. HTTP itself defines \texttt{PUT} and \texttt{DELETE} as idempotent.

\subsection{At-least-once execution, exactly-once acceptance}

Truly running something \emph{exactly once} across unreliable machines is
impossible in general. What the platform guarantees instead is:

\begin{itemize}
\item A job may \emph{run} more than once (a zombie may keep training after it
  was replaced). That is \emph{at-least-once execution}.
\item Exactly one attempt's result is ever \emph{accepted and promoted}. That
  is \emph{exactly-once acceptance}, enforced by fencing and a single
  \texttt{accepted\_attempt\_id} column on the job.
\end{itemize}

Message queues make the same trade: Kafka and SQS deliver at least once, and
consumers make processing idempotent to get an exactly-once \emph{effect}.

\subsection{State machines}

An attempt moves through named states, and only some moves are legal:

\begin{verbatim}
 leased -> starting -> running -> uploading -> completed
    \          \          \          \
     +----------+----------+----------+--> failed | cancelling -> cancelled
                       any active state --> unknown (lease expired)
                       unknown --> back to running (renewed in time) | failed
\end{verbatim}

\begin{lstlisting}[language=Python]
# control/app/enums.py
FORWARD_TRANSITIONS = {
    "leased":    {"starting", "running", "uploading", "failed",
                  "cancelling", "unknown"},
    "running":   {"uploading", "failed", "cancelling", "unknown"},
    "uploading": {"completed", "failed", "cancelling", "unknown"},
    ...
}
\end{lstlisting}

Writing transitions down as data means an impossible report, such as
\texttt{completed} going back to \texttt{running}, becomes a clear
\texttt{409 illegal\_transition} instead of silently corrupting state.

\paragraph{Other use cases.} TCP connections, order fulfilment, CI pipelines,
Kubernetes pod phases and AWS Step Functions are all explicit state machines.
